\section*{الفصل \ref{ch:g3:animals-through-seasons} --- الحيوانات عبر فصول السنة}

\begin{solution}{exo:g3:animals-through-seasons:1}
البرد، وهو يجعل البقاء دافئًا يكلّف \omterm{prop:g3:balanced-diet:jobs}{طاقة} أكثر --- في الوقت نفسه
الذي يشحّ فيه الغذاء، وهو مصدر تلك \omterm{prop:g3:balanced-diet:jobs}{الطاقة}.
\end{solution}

\begin{solution}{exo:g3:animals-through-seasons:2}
\omterm{def:g3:animals-through-seasons:migration}{الهجرة}: السنونو واللقلق (أو الكركي). \omterm{def:g3:animals-through-seasons:hibernation}{والسبات}: القنفذ والغرير
الجبلي (أو الزغبة والخفاش). والبقاء نشطًا: الثعلب والسنجاب (أو
أبو الحنّاء والأيّل).
\end{solution}

\begin{solution}{exo:g3:animals-through-seasons:3}
يبرد جسمه، ويبطؤ قلبه وتنفسه حتى يكادا يقفان، وينام عميقًا في
مأواه. ويعيش على الشحم الذي ادّخره في الخريف.
\end{solution}

\begin{solution}{exo:g3:animals-through-seasons:4}
السنونو يأكل \omterm{prop:g3:sorting-living-things:groups}{الحشرات} الطائرة يمسكها في الجو --- ولا يوجد منها هنا
شيء البتة في الشتاء. وإذ زال مورد غذائه تمامًا، وجب عليه أن يذهب
حيث ما زالت \omterm{prop:g3:sorting-living-things:groups}{الحشرات} تطير: البلاد الأدفأ.
\end{solution}

\begin{solution}{exo:g3:animals-through-seasons:5}
\omterm{def:g3:animals-through-seasons:active}{فروة شتوية} أكثف (كالثعلب والأيّل)؛ \omterm{def:g3:animals-through-seasons:active}{ومخازن غذاء} تُجمع في الخريف
(كجوز السنجاب المدفون)؛ ومائدة مبدَّلة (كالشحرور ينتقل من الديدان
إلى التوت)؛ وبحث أشقّ (كأبي الحنّاء يجوب باحثًا عن آخر التوت).
\end{solution}

\begin{solution}{exo:g3:animals-through-seasons:6}
\omterm{def:g3:animals-through-seasons:migration}{المهاجر} يخاطر برحلة منهكة؛ والسابت عليه أن يجمع كل شحم شتائه في
الخريف وإلا لم يستيقظ في الربيع؛ \omterm{def:g1:animals-around-us:animal}{والحيوان} النشط يقامر على إيجاد
الغذاء عبر أجوع الشهور.
\end{solution}

\begin{solution}{exo:g3:animals-through-seasons:7}
\omterm{def:g3:animals-through-seasons:hibernation}{السبات} --- فالأكل بنهم يبني شحم الشتاء، والحفرة المبطنة \omterm{def:g1:plants-around-us:parts}{بالأوراق}
هي عش الشتاء. والخطوة الرابعة: مراجعة الخريف (فالتسمين وحفر العش
يدلان على \omterm{def:g3:animals-through-seasons:hibernation}{السبات}).
\end{solution}

\begin{solution}{exo:g3:animals-through-seasons:8}
الجواب مفتوح. ولا سنونوات في الشتاء --- فهي في إفريقيا. \omterm{prop:g3:sorting-living-things:groups}{والطيور}
التي تُرى فعلًا (أبو الحنّاء، والقرقف، والشحارير، والعصافير،
والغربان) هي أصحاب خطة البقاء \omterm{def:g3:animals-through-seasons:active}{نشطة}.
\end{solution}

\begin{solution}{exo:g3:animals-through-seasons:9}
(أ) غذاء السنجاب ممكن الإيجاد في الشتاء --- وخاصة المدفون في
الخريف --- فهو إذن نشط، يعيش على مخازنه. (ب) ضفادع الكركي وحشراته
تزول من المروج الرطبة المتجمدة --- فالغذاء زال، وهو إذن \omterm{def:g3:animals-through-seasons:migration}{مهاجر}
(والكراكي تطير جنوبًا فعلًا في أسراب كبيرة على هيئة حرف V).
\end{solution}

\begin{solution}{exo:g3:animals-through-seasons:10}
السابت يدفع ثمن الشتاء سلفًا، بشحم الخريف؛ والقنفذ النحيل لا شيء
عنده يحرقه عبر شهور نومه الباردة فلن يستيقظ. ويأخذ المنقذون قنافذ
الخريف النحيلة ويطعمونها في الداخل حتى تبلغ وزنًا يتيح لها \omterm{def:g3:animals-through-seasons:hibernation}{السبات}
بأمان --- أو يبقونها دافئة مطعَمة طول الشتاء.
\end{solution}
